\documentclass[10pt,twocolumn]{article}

\usepackage[a4paper,margin=0.68in]{geometry}
\usepackage{graphicx}
\usepackage{booktabs}
\usepackage{tabularx}
\usepackage{array}
\usepackage{caption}
\usepackage{hyperref}
\usepackage{titlesec}
\usepackage{microtype}
\usepackage{xurl}
\usepackage{amssymb}
\usepackage{float}

\setlength{\columnsep}{0.26in}
\setlength{\parindent}{0pt}
\setlength{\parskip}{2.5pt}

\graphicspath{{figures/}}

\hypersetup{
    colorlinks=true,
    urlcolor=blue,
    linkcolor=black,
    citecolor=black
}

\captionsetup{
    font=footnotesize,
    labelfont=bf
}

\titlespacing*{\section}{0pt}{5pt}{1.5pt}

\titleformat{\section}
{\normalfont\large\bfseries}
{\thesection.}{0.45em}{}

\begin{document}

\twocolumn[{%
\centering
{\LARGE\bfseries MIT 805 Big Data Semester Project - Part 1}\par
\vspace{2pt}

{\large NYC TLC High-Volume For-Hire Vehicle Trip Records}\par
\vspace{5pt}

{\normalsize
Group [10]\
[Fateh Cheballah Student No. 14161631]\quad
[Mpendulo Sibeko Student No.]
}\par

\vspace{2pt}

{\small GitHub: \url{https://github.com/u14161631-debug/MIT805-Big-Data-Project.git}}\par

\vspace{5pt}
\rule{\textwidth}{0.4pt}
\vspace{6pt}
}]


% ============================================================
\section{Dataset source and suitability}

This project uses the New York City Taxi and Limousine Commission
(NYC TLC) High-Volume For-Hire Vehicle (HVFHV) Trip Record Data.
The records are reported by licensed high-volume for-hire vehicle
bases and published by the TLC as monthly Parquet files
\cite{tlc,dict}.

The dataset is publicly available through the NYC TLC and NYC Open
Data platforms and is suitable for academic use under the NYC Open
Data terms \cite{opendata}. The raw data files are not stored in the
GitHub repository because of their size. Instead, the repository
provides the official source and the code used to identify and access
the required files.


% ============================================================
\section{Dataset characteristics and scale}

The HVFHV dataset contains temporal, spatial, categorical and financial
information. The available fields include request, on-scene, pickup and
drop-off timestamps, pickup and drop-off location IDs, trip distance,
trip duration, fare components, driver pay, operator codes and service
flags.

File sizes were measured directly from the TLC server using HTTP HEAD
requests. The complete available HVFHV corpus contains 89 monthly files
from February 2019 to June 2026 and totals 37.58 GB.

The project working dataset contains 26 recent files totalling 12.00 GB,
while the processing dataset contains 7 files totalling 3.36 GB. These
sizes satisfy the assignment requirements of 25--40 GB for the raw
dataset, at least 12 GB for the working dataset and at least 3 GB for
the processing dataset.

Part 1 exploratory analysis uses the June 2026 file in Polars. This
monthly file contains 20,775,868 records and 25 columns. The larger
3.36 GB multi-file processing dataset is reserved for distributed
PySpark analysis in Part 2.


% ============================================================
\section{Data quality and preparation}

The June 2026 file was kept unchanged and a separate analytical
DataFrame was created for cleaning and EDA. Initial checks identified
4,844 zero-distance trips, 2,514 trips above 100 miles, 6,750 negative
base fares, 1,306 fares above \$500 and 2 non-positive trip durations.
No exact duplicate rows were found.

The main missing field was \texttt{originating\_base\_num}, which was
absent in 26.45\% of records. The main trip timestamp fields were
complete in the June 2026 sample.

Cleaning retained trips with duration in $(0,1440]$ minutes, distance
in $(0,100]$ miles and base passenger fare in $[0,500]$. The combined
filters excluded 14,421 records, equivalent to 0.07\% of the sample,
leaving 20,761,447 trips for analysis.

These thresholds were used as plausibility filters for the EDA rather
than as a claim that every excluded record was necessarily incorrect.
The dataset also has limitations that cleaning cannot address. HVFHV
represents app-based high-volume for-hire activity rather than all NYC
transport, and the one-month EDA cannot capture seasonal patterns.
Passenger count is also not available, so occupancy cannot be analysed.


% ============================================================
\section{Exploratory data analysis}

The cleaned June 2026 sample had a mean trip distance of 4.98 miles
and a median of 2.91 miles. Mean trip duration was 20.36 minutes,
with a median of 16.18 minutes. The mean base passenger fare was
\$28.75 compared with a median of \$20.13.

Trip distance and base passenger fare showed a strong positive
relationship. The correlation was 0.8596 before cleaning and 0.8591
after cleaning, suggesting that the relatively small number of excluded
records had little effect on the overall relationship.

Clear temporal differences were also observed. Saturday recorded the
highest trip count at 3,215,911, compared with 2,695,792 on Wednesday.
Trip characteristics also changed across the day. Trips around 05:00
averaged 7.13 miles, while trips at 15:00 averaged 4.70 miles and
23.97 minutes.

Spatial differences were also visible. LaGuardia Airport was the
largest pickup zone with 421,225 trips, followed by JFK Airport with
356,638. Borough-level pickup totals were approximately 7.14 million
in Manhattan, 5.72 million in Brooklyn, 4.77 million in Queens,
2.79 million in the Bronx and 0.35 million in Staten Island.

The operator field showed that Uber accounted for 73.43\% of trips in
the June 2026 sample, while Lyft accounted for 26.57\%. Selected
figures are included in the Appendix.


% ============================================================
\section{The 7 Vs of Big Data}

\textbf{Volume.}
The available HVFHV corpus contains 89 monthly Parquet files totalling
37.58 GB. The June 2026 file alone contains 20.78 million trip records.
This volume supports the use of a 12.00 GB working dataset and a
3.36 GB processing dataset for distributed analysis.

\textbf{Velocity.}
HVFHV trips are generated continuously through ride-hailing operations.
The June 2026 sample contains roughly 692,500 trips per day on average
and includes request, on-scene, pickup and drop-off timestamps. TLC
publishes the trip files monthly, so the operational data is generated
continuously but released for analysis in batches.

\textbf{Variety.}
The dataset contains 25 columns covering several data types. These
include timestamps, location IDs, trip distance and duration, fare
components, driver pay, operator codes and service flags. This variety
supports temporal, spatial, financial and operator-level analysis.

\textbf{Veracity.}
The quality assessment identified zero-distance trips, negative fares,
extreme trip distances and two non-positive trip durations. After
applying the cleaning rules, 14,421 records, or 0.07\%, were removed.
This shows why basic validation is required before analysing the data.

\textbf{Variability.}
Trip behaviour changes across both time and location. Saturday recorded
3.22 million trips compared with 2.70 million on Wednesday. Average trip
distance also varied by hour, reaching 7.13 miles at 05:00 compared
with 4.70 miles at 15:00. Pickup activity also differed across boroughs
and taxi zones.

\textbf{Visualization.}
A monthly file containing more than 20 million records is difficult to
interpret from raw rows alone. Aggregating the data into hourly,
weekday, company, zone and borough views makes the main patterns and
unusual values easier to identify and communicate.

\textbf{Value.}
The dataset can support demand analysis, driver and vehicle positioning,
airport and borough-level planning, operator comparisons and driver-pay
analysis. Over a longer period, the same dataset could also be used to
study changes in mobility behaviour and transport policy.


% ============================================================
\section{Business and societal value}

The HVFHV dataset can support operational decisions such as identifying
periods of high demand, positioning drivers or vehicles and understanding
airport and borough-level travel activity. Spatial analysis can help show
where demand is concentrated, while temporal analysis can identify when
demand is highest.

Financial and operator fields also allow comparison of trip
characteristics and driver pay across the main HVFHV operators. Over a
longer period, the dataset could support studies of changes in mobility
behaviour and the effects of transport policies.

There are several limitations. HVFHV represents only one part of the
New York City transport system, so the findings cannot be generalised to
all travel in the city. The records are de-identified, which prevents
longitudinal analysis of individual riders or drivers. Locations are
represented by taxi-zone IDs rather than exact GPS coordinates, and some
fields contain missing or unusual values. The June 2026 EDA also
represents only one month.


% ============================================================
\begin{thebibliography}{9}
\footnotesize

\bibitem{tlc}
NYC Taxi \& Limousine Commission.
\emph{TLC Trip Record Data}.
\url{https://www.nyc.gov/site/tlc/about/tlc-trip-record-data.page}

\bibitem{dict}
NYC Taxi \& Limousine Commission.
\emph{High Volume FHV Trip Records Data Dictionary}.
\url{https://www.nyc.gov/assets/tlc/downloads/pdf/data_dictionary_trip_records_hvfhs.pdf}

\bibitem{opendata}
NYC Open Data.
\emph{Frequently Asked Questions}.
\url{https://opendata.cityofnewyork.us/faq/}

\end{thebibliography}


% ============================================================
% APPENDIX
% ============================================================

\clearpage
\onecolumn

\appendix
\section*{Appendix}

\setcounter{table}{0}
\renewcommand{\thetable}{A\arabic{table}}

\setcounter{figure}{0}
\renewcommand{\thefigure}{A\arabic{figure}}


% ------------------------------------------------------------
% Dataset scale table
% ------------------------------------------------------------

\begin{table}[H]
\centering
\caption{Dataset scales used in the project.}
\label{tab:scale}

\begin{tabular}{lrrll}
\toprule
Dataset & Files & Size (GB) & Period & Requirement \\
\midrule
Raw & 89 & 37.58 & Feb 2019--Jun 2026 & 25--40 GB \checkmark \\
Working & 26 & 12.00 & May 2024--Jun 2026 & $\geq$12 GB \checkmark \\
Processing & 7 & 3.36 & Dec 2025--Jun 2026 & $\geq$3 GB \checkmark \\
\bottomrule
\end{tabular}

\end{table}


% ------------------------------------------------------------
% Cleaning summary
% ------------------------------------------------------------

\begin{table}[H]
\centering
\caption{Cleaning summary for the June 2026 EDA sample.}
\label{tab:cleaning}

\begin{tabular}{lrr}
\toprule
Metric & Raw data & Clean data \\
\midrule
Number of records & 20,775,868 & 20,761,447 \\
Zero-distance trips & 4,844 & 0 \\
Negative-fare trips & 6,750 & 0 \\
Non-positive-duration trips & 2 & 0 \\
\bottomrule
\end{tabular}

\end{table}


% ------------------------------------------------------------
% Figure A1 and A2
% ------------------------------------------------------------

\begin{figure}[H]
\centering

\begin{minipage}[t]{0.48\textwidth}
\centering
\includegraphics[width=\linewidth]{hvfhv_size_by_year.png}
\caption{HVFHV dataset size by year.}
\label{fig:size}
\end{minipage}
\hfill
\begin{minipage}[t]{0.48\textwidth}
\centering
\includegraphics[width=\linewidth]{distance_vs_fare_clean.png}
\caption{Trip distance and base passenger fare.}
\label{fig:scatter}
\end{minipage}

\end{figure}


% ------------------------------------------------------------
% Figure A3 and A4
% ------------------------------------------------------------

\begin{figure}[H]
\centering

\begin{minipage}[t]{0.48\textwidth}
\centering
\includegraphics[width=\linewidth]{trips_by_hour.png}
\caption{Trips by pickup hour.}
\label{fig:hour}
\end{minipage}
\hfill
\begin{minipage}[t]{0.48\textwidth}
\centering
\includegraphics[width=\linewidth]{trips_by_day.png}
\caption{Trips by day of week.}
\label{fig:day}
\end{minipage}

\end{figure}


% ------------------------------------------------------------
% Figure A5 and A6
% ------------------------------------------------------------

\begin{figure}[H]
\centering

\begin{minipage}[t]{0.48\textwidth}
\centering
\includegraphics[width=\linewidth]{hourly_characteristics.png}
\caption{Average trip characteristics by hour.}
\label{fig:hourly}
\end{minipage}
\hfill
\begin{minipage}[t]{0.48\textwidth}
\centering
\includegraphics[width=\linewidth]{trips_by_company.png}
\caption{Trips by company.}
\label{fig:company}
\end{minipage}

\end{figure}


% ------------------------------------------------------------
% Figure A7 and A8
% ------------------------------------------------------------

\begin{figure}[H]
\centering

\begin{minipage}[t]{0.48\textwidth}
\centering
\includegraphics[width=\linewidth]{top_pickup_zones.png}
\caption{Top 15 pickup zones.}
\label{fig:zones}
\end{minipage}
\hfill
\begin{minipage}[t]{0.48\textwidth}
\centering
\includegraphics[width=\linewidth]{pickups_by_borough.png}
\caption{Pickups by borough.}
\label{fig:borough}
\end{minipage}

\end{figure}

\end{document}